% Part: first-order-logic
% Chapter: syntax-and-semantics
% Section: unique-readability

\documentclass[../../../include/open-logic-section]{subfiles}

\begin{document}

\olfileid{fol}{syn}{unq}

\olsection{Iedere formule heeft maar één opbouw}

\begin{explain}
Door onze definitie heeft iedere !!{formula} \emph{maar één lezing}:
volgens de vormregels voor !!{formula}s kan zij in wezen maar op één
manier zijn opgebouwd. Daardoor kan zij ook maar op één manier worden
``geïnterpreteerd''. Zonder die eigenschap zouden sommige !!{formula}s
meer dan één lezing of interpretatie hebben. Dat willen we voorkomen.
Nog belangrijker is dat de meeste definities en bewijzen hierna anders
niet zouden werken.

Een slecht voorbeeld laat het probleem het duidelijkst zien. Stel dat
onze regels voor !!{formula}s meer dan één opbouw toelieten. We zouden
bijvoorbeeld de haakjes rond verbindingstekens kunnen vergeten en deze
regel kunnen toelaten:
\begin{quote}
Als $!A$ en $!B$ !!{formula}s zijn, dan is $!A \lif !B$ dat ook.
\end{quote}
Begin met een atomaire formule $!D$. Met deze regel kunnen we dan $!D
\lif !D$ maken. Met die formule en $!D$ kunnen we vervolgens $!D \lif
!D \lif !D$ maken. Dat kan op twee manieren:
\begin{enumerate}
\item We nemen $!D$ als $!A$ en $!D \lif !D$ als $!B$.
\item We nemen $!A$ als $!D \lif !D$ en $!B$ als~$!D$.  
\end{enumerate}
Daarom kunnen we de !!{formula}~$!D \lif !D \lif !D$ ook op twee
manieren ``lezen''. Zij heeft de vorm $!B \lif !C$ met $!B$ gelijk aan
$!D$ en $!C$ gelijk aan $!D \lif !D$. Maar zij heeft \emph{ook} de vorm
$!B \lif !C$ met $!B$ gelijk aan $!D \lif !D$ en $!C$ gelijk aan~$!D$.

Dan werken onze definities niet altijd. Neem de !!{main operator} van
een formule. Bij een formule van de vorm $!B \lif !C$ noemen we het
aangegeven voorkomen van~$\lif$ de !!{main operator}. De formule $!D
\lif !D \lif !D$ past hierboven op twee manieren bij $!B \lif !C$.
Bij de ene manier wordt het eerste voorkomen van $\lif$ de !!{main
operator}; bij de andere manier wordt het tweede voorkomen dat. Maar
de !!{main operator} moet een \emph{functie} van de !!{formula} zijn.
Dat wil zeggen: iedere !!{formula} moet precies één voorkomen van een
!!{main operator} hebben.
\end{explain}

\begin{lem}
In iedere !!{formula}~$!A$ staan evenveel linker- als rechterhaakjes.
\end{lem}

\begin{proof}
We bewijzen dit door te volgen hoe $!A$ is opgebouwd. Dit heet
inductie over de opbouw. We moeten twee dingen doen. Eerst controleren
we het voor alle atomaire formules; dat is de basis van de inductie.
Daarna laten we zien dat iedere vormregel de eigenschap bewaart: als
de formules waarmee we beginnen evenveel linker- als rechterhaakjes
hebben, dan geldt dat ook voor de nieuwe formule die we ermee maken.

We schrijven $l(!A)$ voor het aantal linkerhaakjes en $r(!A)$ voor het
aantal rechterhaakjes in~$!A$. Op dezelfde manier zijn $l(t)$ en $r(t)$
de aantallen linker- en rechterhaakjes in een term~$t$.

\begin{prob}
    Bewijs dat voor iedere term~$t$ geldt: $l(t) = r(t)$.
\end{prob}

\begin{enumerate}
\tagitem{prvFalse}{\indcase{!A}{\lfalse}{In $\indfrm$ staan $0$
    linker- en $0$ rechterhaakjes.}}{}

\tagitem{prvTrue}{\indcase{!A}{\ltrue}{In $\indfrm$ staan $0$
    linker- en $0$ rechterhaakjes.}}{}

\item \indcase{!A}{\Atom{R}{t_1,\dots,t_n}}{$l(\indfrm) = 1 + l(t_1) +
  \dots + l(t_n) = 1 + r(t_1) + \dots + r(t_n) = r(\indfrm)$. Hier
  gebruiken we de oefening hierboven: $l(t) = r(t)$ voor iedere
  term~$t$.}

\item \indcase{!A}{\eq[t_1][t_2]}{$l(\indfrm) = l(t_1) + l(t_2) =
  r(t_1) + r(t_2) = r(\indfrm)$.}

\tagitem{prvNot}{\indcase{!A}{\lnot !B}{Volgens de inductieaanname
    geldt $l(!B) = r(!B)$. Daarom geldt $l(\indfrm) = l(!B) = r(!B) =
    r(\indfrm)$.}}{}

\item \indcase{!A}{(!B \ast !C)}{Volgens de inductieaanname geldt $l(!B) =
  r(!B)$ en $l(!C) = r(!C)$. Daarom geldt $l(\indfrm) = 1 + l(!B) + l(!C) =
  1 + r(!B) + r(!C) = r(\indfrm)$.}

\tagitem{prvAll}{\indcase{!A}{\lforall[x][!B]}{Volgens de
    inductieaanname geldt $l(!B) = r(!B)$. Daarom geldt $l(\indfrm) = l(!B) = r(!B) =
    r(\indfrm)$.}}{}

% Toon het geval voor \lexists als prvEx geldt en prvAll niet
\tagitem{prvAll,notprvEx}{}{\indcase{!A}{\lexists[x][!B]}{Volgens de
    inductieaanname geldt $l(!B) = r(!B)$. Daarom geldt $l(\indfrm) = l(!B) = r(!B) =
    r(\indfrm)$.}}

% Zeg alleen 'Hetzelfde' als zowel prvEx als prvAll gelden
\tagitem{notprvAll,notprvEx}{}{\indcase{!A}{\lexists[x][!B]}{Hetzelfde.}}
\end{enumerate}
\end{proof}

\begin{defn}[Echt beginstuk (echte prefix)]
Een rij tekens $!B$ is een \emph{echt beginstuk}, ook wel een echte
prefix, van een rij tekens~$!A$ als je na $!B$ nog een niet-lege rij
tekens kunt zetten en zo precies~$!A$ krijgt.
\end{defn}

\begin{lem}\ollabel{lem:no-prefix}
Als $!A$ !!a{formula} is en $!B$ een echt beginstuk van $!A$ is, dan
is $!B$ geen !!a{formula}.
\end{lem}

\begin{proof}
Oefening.
\end{proof}

\begin{prob}
Bewijs \olref[fol][syn][unq]{lem:no-prefix}.
\end{prob}

\begin{prop}
\ollabel{prop:unique-atomic}
Als $!A$ een atomaire !!{formula} is, geldt precies één van de volgende
mogelijkheden.
\begin{enumerate}
\tagitem{prvFalse}{$!A \ident \lfalse$.}{}
\tagitem{prvTrue}{$!A \ident \ltrue$.}{}
\item $!A \ident \Atom{R}{t_1,\dots,t_n}$, waarbij $R$ een $n$-plaatsig
  !!{predicate} is en $t_1$, \dots, $t_n$ termen zijn. De symbolen $R$,
  $t_1$, \dots, $t_n$ liggen allemaal precies vast.
\item $!A \ident \eq[t_1][t_2]$, waarbij er voor $t_1$ en $t_2$ maar
  één keuze mogelijk is.
\end{enumerate}
\end{prop}

\begin{proof}
Oefening.
\end{proof}

\begin{prob}
Bewijs \olref[fol][syn][unq]{prop:unique-atomic}. (Aanwijzing: formuleer
en bewijs voor termen een versie van \olref[fol][syn][unq]{lem:no-prefix}.)
\end{prob}

\begin{prop}[Eenduidige opbouw]
Iedere !!{formula} valt onder precies één van de volgende mogelijkheden.
\begin{enumerate}
\item $!A$ is atomair.

\tagitem{prvNot}{$!A$ heeft de vorm $\lnot !B$.}{}

\tagitem{prvAnd}{$!A$ heeft de vorm $(!B \land !C)$.}{}

\tagitem{prvOr}{$!A$ heeft de vorm $(!B \lor !C)$.}{}

\tagitem{prvIf}{$!A$ heeft de vorm $(!B \lif !C)$.}{}

\tagitem{prvIff}{$!A$ heeft de vorm $(!B \liff !C)$.}{}

\tagitem{prvAll}{$!A$ heeft de vorm $\lforall[x][!B]$.}{}

\tagitem{prvEx}{$!A$ heeft de vorm $\lexists[x][!B]$.}{}
\end{enumerate}
In elk geval ligt $!B$, of liggen $!B$ en $!C$, precies vast. Er
zijn dus bijvoorbeeld geen twee verschillende paren $!B$, $!C$ en
$!B'$, $!C'$ waarvoor $!A$ tegelijk de vorm
\iftag{prvIf}{$(!B \lif !C)$ en $(!B' \lif !C')$ heeft.}{
    \iftag{prvOr}{$(!B \lor !C)$ en $(!B' \lor !C')$ heeft.}{
      \iftag{prvAnd}{$(!B \land !C)$ en $(!B' \land !C')$ heeft.}{}}}
\end{prop}

\begin{proof}
Volgens de vormregels begint iedere niet-atomaire !!{formula} met een
openingshaakje~(, \iftag{prvNot}{$\lnot$,}{} of een teken voor
``voor alle'' of ``er bestaat''. Begint een
!!{formula} met een van de volgende symbolen, dan moet zij juist
atomair zijn: !!a{predicate}, !!a{function}, !!a{constant}\iftag{prvFalse}{,
$\lfalse$}{}\iftag{prvTrue}{, $\ltrue$}{}.

We hoeven daarom alleen nog het volgende geval te behandelen. Stel dat
$!A$ de vorm $(!B \ast !C)$ én de vorm $(!B' \mathbin{\ast'} !C')$
heeft. Dan moeten we bewijzen dat $!B \ident !B'$, $!C \ident !C'$ en
$\ast = {\ast'}$.

Neem dus aan dat zowel $!A \ident (!B \ast !C)$ als $!A \ident (!B'
\mathbin{\ast'} !C')$ geldt. Nu zijn er twee mogelijkheden: $!B \ident
!B'$ geldt wel, of niet. Als het wel geldt, volgen $\ast = {\ast'}$ en
$!C \ident !C'$ meteen. Die tekens zijn namelijk even lange stukken
van $!A$ die op dezelfde plaats beginnen. In het andere geval geldt
$!B \not\ident !B'$. De stukken $!B$ en $!B'$ beginnen allebei op
dezelfde plaats in $!A$. Daarom moet het ene een echt beginstuk van het
andere zijn. Volgens \olref{lem:no-prefix} kan dat niet.
\end{proof}
\end{document}
